% =============================================================================
% dft.tex — Draft: Trustless Coordination: How DAO Governance Achieves
%           Consensus Without Hierarchy
%
% Target: AOM/SMS Conference — Technology and Innovation Management (TIM) track
% Format: IEEE conference (IEEEtran.cls)
% Date:   2026-03-15
% Status: Working draft — all statistics from annotated_records.json (N=5,421
%         total), structural_metrics.csv, core_contributors.csv.
%         Institution data: R07 manual investigation (109 authors verified).
%
% NOTE TO AUTHOR: Sections marked [AUTHOR] require your judgment or additions.
%   Numbers in this draft come directly from the data pipeline; please verify
%   against your own reading of key records before submission.
% =============================================================================

\documentclass[conference]{IEEEtran}
\IEEEoverridecommandlockouts

\usepackage{cite}
\usepackage{amsmath,amssymb,amsfonts}
\usepackage{graphicx}
\usepackage{textcomp}
\usepackage{xcolor}
\usepackage{booktabs}
\usepackage{array}
\usepackage{url}

% Tighten table column spacing
\setlength{\tabcolsep}{4pt}

\def\BibTeX{{\rm B\kern-.05em{\sc i\kern-.025em b}\kern-.08em
    T\kern-.1667em\lower.7ex\hbox{E}\kern-.125emX}}

\begin{document}

% =============================================================================
\title{Trustless Coordination: How Permissionless DAO Governance Achieves
       Consensus Under Conditions That Transaction Cost Theory Predicts
       Require Hierarchy}
% =============================================================================

\author{
  \IEEEauthorblockN{[Author Name]}
  \IEEEauthorblockA{
    \textit{[Department]} \\
    \textit{[Institution]} \\{}
    {[City, Country]} \\{}
    {[email]}
  }
}

\maketitle

% =============================================================================
\begin{abstract}
% =============================================================================
Transaction cost economics (TCE) predicts that high asset specificity and
behavioral uncertainty jointly favor hierarchical governance over markets or
peer-to-peer coordination. Yet permissionless decentralized autonomous
organization (DAO) governance increasingly produces binding technical standards
in exactly these conditions. Drawing on TCE and open-source governance theory,
we conduct a comparative case study of two AI agent protocol governance
processes: ERC-8004 (a Trustless Agents standard proposed through the Ethereum
Improvement Proposal process) and Google A2A (an Agent-to-Agent protocol
governed by a corporate-led technical steering committee). Analysis of 5,421
archival governance records reveals that the DAO process achieved ratified
consensus in 169 days with 69 unique contributors across competing organizations, while
the corporate process remains in ongoing iteration despite 771 participants and
an open-source facade. We identify three mechanisms---smart-contract-enforced
commitment, pseudonymous reputation externalization, and permissionless
competitive entry---that jointly substitute for hierarchical authority. Our
findings contribute to new forms of organizing theory by specifying the
operational conditions under which blockchain governance outperforms hierarchy,
and extend TCE by identifying code-enforced credible commitment as a
fourth governance mechanism alongside markets, hybrids, and hierarchy.
\end{abstract}

\begin{IEEEkeywords}
DAO governance, transaction cost economics, technology standardization,
open-source governance, blockchain, AI protocol, credible commitment
\end{IEEEkeywords}


% =============================================================================
\section{Introduction}
% =============================================================================

How does collective action among competing firms produce binding technical
standards without contracts, authority, or equity alignment? This question is
not new to organizational theory---open standards bodies have operated for
decades---but it takes on fresh urgency as decentralized autonomous organizations
(DAOs) begin producing governance outcomes that challenge the canonical
predictions of transaction cost economics~\cite{williamson1985}.

The puzzle is sharpest when the conditions are most demanding. Williamson's
comparative governance framework identifies two dimensions that make hierarchical
authority efficient: \textit{asset specificity} (investments that lose value
outside the relationship) and \textit{behavioral uncertainty} (inability to
specify contingencies ex ante)~\cite{williamson1991}. When both are high, TCE
predicts that hierarchical governance outperforms markets and peer-to-peer
coordination because only authority can prevent opportunistic hold-up. Yet
beginning in 2025, the Ethereum developer community coordinated a multi-firm
AI agent interoperability protocol---ERC-8004---through a permissionless process
with no employment authority, no equity stakes, and no contracts, reaching
binding consensus in 169 days. Both asset specificity and behavioral uncertainty
were high: the standard required coordinated implementation across wallet
providers, agent frameworks, and decentralized applications, in a novel domain
with no prior art.

This paper documents the puzzle and resolves it. We compare ERC-8004's DAO
governance with Google A2A, a contemporaneous AI agent protocol governed by a
corporate-led Technical Steering Committee (TSC) with formal membership and
binding votes. Both cases involve AI protocol standardization; they differ
systematically in governance mode. Analysis of 5,421 archival governance
records---scraped from Ethereum Magicians forum, GitHub repositories, and
GitHub Discussions---reveals three mechanisms by which the DAO process achieves
coordination without hierarchical authority.

Our contributions are threefold. First, we document an empirical regularity
that the conventional TCE prediction struggles to accommodate: permissionless
governance achieves binding consensus faster than a corporate hierarchy
under theoretically hierarchy-inducing conditions. Second, we identify the
operational mechanisms: smart contracts as credible commitment devices,
pseudonymous reputation as governance currency, and permissionless entry
as anti-capture discipline. Third, we contribute to the new forms of organizing
literature~\cite{puranam2014} by specifying boundary conditions under which
blockchain governance is organizationally efficient.


% =============================================================================
\section{Theoretical Background}
% =============================================================================

\subsection{Transaction Cost Economics and Governance Mode}

Transaction cost economics provides the canonical framework for predicting
governance mode selection. Williamson~\cite{williamson1985} identifies three
key transaction characteristics: asset specificity, uncertainty, and frequency.
When asset specificity and uncertainty are high, the \textit{hold-up problem}
makes market governance hazardous: ex post opportunism by counterparties who
know that specific investments are sunk. Hierarchical governance solves
this by bringing transactions inside the authority relation, where managerial
fiat can resolve disputes without the threat of defection~\cite{williamson1991}.

Between market and hierarchy lies the hybrid, including alliances, joint
ventures, and standards bodies~\cite{gulati1998}. But even hybrid governance
requires some formal structure---membership criteria, voting rules, intellectual
property agreements---to deter opportunism. TCE offers no clean governance
category for a permissionless process that imposes no membership requirements,
issues no contracts, and enforces standards through code rather than legal
sanction. The DAO governance mode is thus theoretically anomalous under the
standard framework.

\subsection{Open-Source and Commons Governance}

The open-source software (OSS) literature provides partial theoretical
resources. O'Mahony and Ferraro~\cite{omahony2007} document how OSS communities
develop governance structures endogenously, with meritocratic authority emerging
from contribution history. Shah~\cite{shah2006} identifies how reputation
mechanisms sustain hybrid governance without employment hierarchy. Simcoe's
\cite{simcoe2012} analysis of standard-setting organizations (SSOs) is the
closest analog: SSO committees achieve consensus governance for shared technology
platforms, but they do so through formal membership gates that the EIP process
lacks.

The fork threat is central to OSS governance theory. West~\cite{west2003}
argues that the credible option of forking disciplines platform owners even
when the threat is never exercised. In the EIP context, any participant can
propose a competing standard, which imposes competitive discipline on
incumbent actors who might otherwise capture the process.

\subsection{Blockchain Governance as a Novel Organizational Form}

Lumineau et al.~\cite{lumineau2021} identify blockchain governance as a new
way of organizing collaborations, distinct from both markets and hierarchies.
Smart contracts provide what Williamson~\cite{williamson1983} termed credible
commitment: a promise made binding through structure rather than words. Unlike
contract law, which requires enforcement by courts, smart-contract commitments
are self-executing---deviation is technically infeasible, not merely
contractually prohibited~\cite{defilippi2018}.

Puranam et al.~\cite{puranam2014} ask what is ``new'' about new forms of
organizing: their answer is that novel organizational forms address
coordination and cooperation problems differently, not necessarily better.
We build on this: the EIP process addresses the hold-up problem through
code-enforced commitment rather than authority, and addresses opportunism
through reputation rather than supervision. Whether this constitutes a
new governance mode or an extension of the hybrid requires specifying the
mechanisms by which it achieves coordination---the contribution of this paper.

\subsection{Research Questions}

This paper addresses two research questions:

\begin{description}
  \item[RQ1] Under conditions that TCE predicts require hierarchical governance,
             how does permissionless DAO governance achieve binding consensus?
  \item[RQ2] What mechanisms substitute for hierarchical authority in DAO
             governance, and how do they compare to corporate governance of
             technology standards?
\end{description}


% =============================================================================
\section{Research Design and Methods}
% =============================================================================

\subsection{Comparative Case Study Design}

We adopt a comparative case study design~\cite{eisenhardt1989, yin2018},
selecting ERC-8004 and Google A2A for theoretical contrast. Both are AI agent
protocol governance processes active in 2025--2026, making the comparison
contemporaneous and domain-matched. They differ systematically on the
theoretically relevant dimension: ERC-8004 operates through permissionless
DAO governance, while Google A2A operates through corporate-led governance
with a formal TSC. This design follows Eisenhardt's~\cite{eisenhardt1989}
logic of selecting cases that are ``polar types'' on the dimension of
theoretical interest.

Neither case was selected for statistical representativeness. Both are
successful protocols by observable outcome (ERC-8004 reached mainnet; A2A
is actively maintained), which introduces survivorship bias that we address
in limitations. The comparison is intended to illuminate mechanism, not to
generalize prevalence.

\subsection{Case Descriptions}

\textit{ERC-8004 (Trustless Agents)} was proposed on August 13, 2025, with
four institutional co-authors listed in the EIP header: Davide Crapis
(\texttt{davide@ethereum.org}, Ethereum Foundation), Marco De Rossi
(\texttt{@MarcoMetaMask}, MetaMask), Jordan Ellis
(\texttt{jordanellis@google.com}, Google), and Erik Reppel
(\texttt{erik.reppel@coinbase.com}, Coinbase). The standard defines a
registry for trustless AI agents on the Ethereum Virtual Machine (EVM).
It was processed through the Ethereum Improvement Proposal (EIP)
pipeline---public deliberation on Ethereum Magicians forum and GitHub---and
merged to mainnet on January 29, 2026. The process required no membership,
no fee, and no organizational affiliation to participate.

Notably, Google and Coinbase appear as formal co-authors yet have zero
public discussion records in the scraped data. Their institutional commitment
was expressed through co-authorship, not through discussion-floor participation.

\textit{Google A2A (Agent-to-Agent)} was first committed on March 25, 2025
and publicly announced on April 9, 2025. The protocol defines an open
standard for AI agent interoperability across vendors. Despite being hosted
as an open-source repository with 771 non-bot contributors, governance authority
rests with a Technical Steering Committee (TSC) whose membership is formally
registered; binding votes are cast through a git-vote mechanism accessible
only to TSC members.

\subsection{Data Collection}

We collected archival data from three sources for each case. For ERC-8004:
(1) 113 posts from the Ethereum Magicians forum thread (topic 25098), scraped
via the Discourse JSON API; and (2) 36 GitHub records from 9 core lifecycle
pull requests (\#1170, \#1244, \#1248, \#1458, \#1462, \#1470, \#1472,
\#1477, \#1488) that directly modified \texttt{ERCS/erc-8004.md}.

For Google A2A: (1) 3,104 issue and issue comment records; (2) 1,955 pull
request and review comment records; (3) 822 GitHub Discussion records scraped
via the GitHub GraphQL API; and (4) 142 records from the two TSC-voted pull
requests (\#831 and \#1206), scraped using the GitHub REST API.

Total records before filtering: 6,172. After removing records with fewer than
20 characters of text (bot messages, CI notifications) and bot authors: 5,421
(ERC-8004: 149; Google A2A: 5,272).
Data collection occurred between November 2025 and March 2026.

\subsection{LLM-Assisted Annotation}

All records were annotated using the MiniMax-M2.5 reasoning model via a
structured prompt (available in the replication archive). Each record received four
categorical labels:
\begin{itemize}
  \item \textbf{Stakeholder institution}: Google, MetaMask, Ethereum
        Foundation, Coinbase, Independent, Unknown
  \item \textbf{Argument type}: Technical, Economic, Governance-Principle,
        Process, Off-topic
  \item \textbf{Stance}: Support, Oppose, Modify, Neutral, Off-topic
  \item \textbf{Consensus signal}: Adopted, Rejected, Pending, N/A
\end{itemize}

The model's reasoning output (\texttt{<think>...</think>} blocks) was
stripped before JSON parsing. Annotation completion rate was 98.7\%
(122 records failed due to text-too-short or JSON parse errors). Institution
labels follow a three-tier provenance cascade: (1) EIP header email
(highest confidence, 2 authors); (2) R07 manual investigation---a systematic
review of GitHub profiles, LinkedIn, and EIP metadata for the top 109
contributors across both cases, yielding 40 institution upgrades---;
and (3) LLM inference (517 remaining authors). The original LLM-inferred label
is preserved in a separate \texttt{institution\_lm} field for all records.

\subsection{Validity Measures}

\noindent\textit{Construct validity:} The openness index (proportion of
non-initiating-institution contributors) operationalizes the theoretical
concept of permissionlessness. Governance escalation events are identified
as records containing formal voting commands (\texttt{/vote}) or editorial
merge decisions (pull request state = merged).

\noindent\textit{Internal validity:} Process tracing links observed patterns
(e.g., high Technical argument proportion) to the theoretical mechanism
(permissionless meritocracy). The two TSC-voted pull requests provide
natural experiments in governance escalation.

\noindent\textit{Limitations on reliability:} Inter-coder reliability
between LLM annotation and human codes has not been computed for the full
sample. Institution labels for 517 of 626 unique authors remain LLM-inferred;
109 were manually verified via R07 investigation, correcting 40 labels
(including material cases such as \texttt{pcarranzav} to The Graph and
\texttt{davidecrapis.eth} to Ethereum Foundation). [AUTHOR: compute Cohen's
$\kappa$ on a random subsample of $N=50$ and report before submission.]


% =============================================================================
\section{Results}
% =============================================================================

\subsection{Speed and Decision Architecture}

Table~\ref{tab:structural} presents the structural governance metrics for
both cases. ERC-8004 reached ratified consensus in 169 days, with a
clear terminus: EIP mainnet inclusion. Google A2A remains in ongoing
iteration with no equivalent endpoint.

\begin{table}[h]
\caption{Structural Governance Metrics}
\label{tab:structural}
\centering
\begin{tabular}{lrr}
\toprule
\textbf{Dimension} & \textbf{ERC-8004} & \textbf{Google A2A} \\
\midrule
Proposal date & 2025-08-13 & 2025-03-25 \\
Consensus achieved & 2026-01-29 & Ongoing \\
Days to consensus & \textbf{169} & — \\
Discussion records & 149 & 5,272 \\
Unique contributors (non-bot) & 69 & 771 \\
Top institution share & MetaMask 12.6\% & Google 25.0\% \\
Independent share & 44.4\% & 44.1\% \\
Formal voting & None & TSC git-vote (2 PRs) \\
\bottomrule
\end{tabular}
\end{table}

The 15-day gap between A2A's first commit (2025-03-25) and public
announcement (2025-04-09) implies significant internal pre-work before
external participation was enabled---a structural asymmetry absent from
the EIP process.

\subsection{Argument Type Distribution}

Table~\ref{tab:arguments} compares argument type distributions across
the two cases.

\begin{table}[h]
\caption{Argument Type Distribution by Case (\%)}
\label{tab:arguments}
\centering
\begin{tabular}{lrr}
\toprule
\textbf{Argument Type} & \textbf{ERC-8004} & \textbf{Google A2A} \\
\midrule
Technical          & \textbf{74.3} & 62.7 \\
Process            & 13.9 & \textbf{28.4} \\
Governance-Principle & 4.9 & 1.7 \\
Off-topic          & 4.2 & 7.0 \\
Economic           & 2.8 & 0.1 \\
\midrule
Total ($N$)        & 149 & 5,272 \\
\bottomrule
\end{tabular}
\end{table}

ERC-8004 deliberation is dominated by Technical arguments (74.3\%),
significantly higher than A2A's 62.7\% ($\chi^2$ test, $p<0.01$ [AUTHOR:
verify]). A2A's Process argument proportion (28.4\% vs.\ 13.9\%) is
the fingerprint of bureaucratic coordination overhead: TSC review workflows,
versioning debates, contribution guidelines, and PR labeling. The EIP
process externalizes process governance through its immutable stage
framework (Draft $\to$ Review $\to$ Last Call $\to$ Final), freeing
participants to focus on technical substance.

The near-absence of Economic arguments in both cases ($\leq 3\%$) is
noteworthy: neither governance process is primarily a political or
economic negotiation. This is consistent with both cases operating in
the technical standards domain, where contribution legitimacy derives
from engineering quality, not negotiating leverage.

\subsection{Stance Distribution and Adversarial Review}

\begin{table}[h]
\caption{Stance Distribution by Case (\%)}
\label{tab:stance}
\centering
\begin{tabular}{lrr}
\toprule
\textbf{Stance} & \textbf{ERC-8004} & \textbf{Google A2A} \\
\midrule
Neutral  & 31.9 & 46.7 \\
Support  & 30.6 & 24.1 \\
Modify   & 29.2 & 21.8 \\
Oppose   & \textbf{6.9} & 2.7 \\
Off-topic & 1.4 & 4.5 \\
\midrule
Total ($N$) & 149 & 5,272 \\
\bottomrule
\end{tabular}
\end{table}

The near-equal three-way split among Neutral (31.9\%), Support (30.6\%),
and Modify (29.2\%) in ERC-8004 is a signature of genuine deliberation.
A rubber-stamp process would show $>70\%$ Support; a captured process
would show concentrated Modify from a small set of powerful actors.

ERC-8004's Oppose rate (6.9\%) is more than double A2A's (2.7\%). The
10 Oppose records in ERC-8004 originated primarily from non-EF contributors---
\texttt{pcarranzav} (16 records, The Graph) with the highest Modify+Oppose rate,
and \texttt{spengrah} (Hats Protocol)---neither of whom holds any procedural
authority in the EIP process. Their objections forced specification revisions
across multiple PR iterations, demonstrating that technical credibility alone
confers influence in the DAO governance model.

A2A's lower Oppose rate may reflect selection effects (strong opponents
exit rather than engage), social inhibition (opposing Google's core
committers carries implicit career risk), or structural displacement
(TSC votes resolve disputes before they appear in the public record
as explicit opposition).

\subsection{Institutional Concentration}

\begin{table}[h]
\caption{Record-Level Institutional Distribution (\%)}
\label{tab:institutions}
\centering
\begin{tabular}{lrr}
\toprule
\textbf{Institution} & \textbf{ERC-8004 (\%)} & \textbf{Google A2A (\%)} \\
\midrule
Independent          & \textbf{44.4} & \textbf{44.1} \\
Unknown              & 13.3 & 15.1 \\
MetaMask             & \textbf{12.6} & — \\
Ethereum Foundation  & 10.4 & — \\
The Graph            & 5.9 & — \\
Google               & — & \textbf{25.0} \\
Microsoft            & — & 9.4 \\
Cisco                & — & 3.8 \\
Other named          & 13.4 & 2.6 \\
\bottomrule
\end{tabular}
\end{table}

The independent contributor share is nearly identical across cases (44.4\%
vs.\ 44.1\%), suggesting that participation breadth is not the differentiating
dimension.

The starkest institutional contrast is the role of Google. In A2A, Google
accounts for 25.0\% of all discussion records, with \texttt{holtskinner}
alone contributing $\approx$560 records ($\approx$10\% of all A2A records).
In ERC-8004, Google is listed as a formal co-author (\texttt{jordanellis@google.com}
in the EIP header) yet has \textit{zero} public discussion records. The same
pattern holds for Coinbase (\texttt{erik.reppel@coinbase.com}, formal co-author,
zero discussion records). Both institutions expressed institutional commitment
through co-authorship, then stepped back from the public deliberative process.
This is the governance complement of the permissionless mechanism: even major
players who endorsed the standard did not leverage their institutional weight
to dominate the specification-writing discourse.

ERC-8004's top institution by discussion volume (MetaMask, 12.6\%) is less
than half of Google's A2A share, and MetaMask's contribution intensity was
episodic---reversed when The Graph's \texttt{pcarranzav} and Hats Protocol's
\texttt{spengrah} raised substantive objections that forced specification
revisions. No equivalent independent-driven reversal appears in the A2A record.

\subsection{Governance Escalation Events}

Two pull requests in the A2A dataset triggered formal TSC votes, representing
\textit{governance escalation events}---situations where normal consensus
mechanisms were insufficient.

\textbf{PR \#831 (tasks/list method, PASSED):} After four months of
discussion (July--November 2025), \texttt{amye} (TSC member) triggered
a \texttt{/vote} after ordinary review stalled on versioning questions.
The vote formalized what was near-consensus; its necessity reveals a gap
in A2A's governance design: without a formal ``approve-and-merge''
endpoint, TSC members lacked a mechanism to close substantive but
slow-moving threads.

\textbf{PR \#1206 (lastUpdateTime field, CLOSED/SUPERSEDED):} A design
dispute between adding a new field vs.\ renaming a filter parameter was
resolved at a private TSC meeting (January 6, 2026), with the decision
then reflected back into GitHub. Discussion migrated to Discord before
resolution, permanently losing the public record. A \texttt{/cancel-vote}
was issued after the TSC call.

This is the most analytically significant governance failure in the
dataset. A substantive technical decision---the naming and semantics
of a data model field---was resolved in a private synchronous call on
a proprietary messaging platform. For ERC-8004, every equivalent
decision appears as a dated, public forum post or PR comment. The
\textit{opacity gradient} between the two cases is highest precisely
when governance is most contested.

No equivalent governance escalation events exist in the ERC-8004 dataset:
the EIP stage structure (Draft $\to$ Review $\to$ Last Call $\to$ Final)
provides a clear endpoint that does not require escalation to a formal vote.


% =============================================================================
\section{Discussion}
% =============================================================================

\subsection{Three Mechanisms of Trustless Coordination}

The data support three mechanisms by which DAO governance achieves
coordination that TCE theory predicts should require hierarchical authority.
We term the overall configuration \textit{trustless coordination}: governance
without hierarchy, sustained by code-enforced commitments, reputational
bonding, and competitive contestability.

\subsubsection{Mechanism 1: Smart Contract as Credible Commitment Device}

Once ERC-8004 reaches the ``Last Call'' stage and is merged to
\texttt{ethereum/ERCs}, the standard becomes part of the canonical Ethereum
specification. Implementation is enforced not by contract or employment
hierarchy but by the EVM's consensus rules---any contract deviating from
the spec will fail verification. This is the blockchain equivalent of
Williamson's credible commitment~\cite{williamson1983}: the cost of defection
is technical incompatibility, not legal sanction.

This mechanism explains the observed intensity of technical review.
ERC-8004's 74.3\% Technical argument proportion reflects participants'
awareness that post-adoption renegotiation is technically infeasible. The
deliberation substitutes for a contractual ex ante specification of
contingencies---exactly what TCE identifies as the solution to incomplete
contracting~\cite{williamson1985}. The difference is that the EIP process
achieves this through code rather than contract.

\subsubsection{Mechanism 2: Pseudonymous Reputation Externalization}

ERC-8004 contributors operate under semi-persistent pseudonyms
(Ethereum addresses, GitHub handles) tied to a public, tamper-proof
contribution history. A contributor who reverses a position without
explanation suffers reputational cost with any participant who inspects
the \texttt{git blame} history. This is reputational bonding without
employment hierarchy---the mechanism Shah~\cite{shah2006} theorized for OSS
governance, amplified by the permanent auditability of blockchain-anchored
identity systems.

The empirical signature is \texttt{pcarranzav}'s consistent Modify/Oppose
stance across 16 records representing a coherent technical position. The
EIP was revised in response to concerns from a contributor from The Graph
(Edge and Node), who holds no procedural authority in the EIP process---no
membership seat, no editorial role. Technical credibility alone conferred
influence. Under corporate hierarchy (A2A), the equivalent mechanism is
employment reputation, which carries career risk that suppresses opposition
from non-employees.

\subsubsection{Mechanism 3: Permissionless Entry as Anti-Capture Mechanism}

The EIP process allows any individual to submit a standard proposal,
comment on any PR, and signal on any EIP without credentialing,
membership dues, or organizational affiliation. This creates an
\textit{unlimited threat of competitive entry}~\cite{west2003}: if ERC-8004
were perceived as captured by MetaMask, independent developers could fork
the EIP and propose a competing standard. The mere possibility disciplines
institutional participants.

The evidence is twofold. First, MetaMask's 12.6\% record share shows high
institutional intensity but not monopoly, and institutional dominance did not
suppress the Oppose voices from The Graph and Hats Protocol that forced
specification revisions. Second---and more revealing---Google and Coinbase
co-authored the EIP (confirmed via EIP header email addresses) yet contributed
zero public discussion records. The permissionless entry threat worked in
reverse: even formal institutional co-authors did not exploit their status
to claim editorial authority over the specification. Google A2A's TSC gate
is the structural opposite: binding authority is limited to a pre-registered
committee, and institutional actors maintain persistent directional control
through contribution volume.

\subsection{Theoretical Implications}

\noindent\textit{For TCE:} The standard governance triad---market, hierarchy,
hybrid---does not have an adequate category for DAO governance. The three
mechanisms identified here collectively perform the function that TCE
assigns to hierarchy: reducing opportunism risk under high asset specificity
and uncertainty. We propose a fourth governance mode, code-enforced
coordination (or ``trustless coordination''), characterized by: (a) self-enforcing
commitment through smart contracts, (b) reputation bonding without employment
authority, and (c) competitive discipline through permissionless entry.
This extends Williamson's framework rather than refuting it.

\noindent\textit{For new forms of organizing:} Puranam et al.~\cite{puranam2014}
argue that new organizational forms are distinguished by how they address
the fundamental problems of coordination (``who does what'') and cooperation
(``who gets what''). The EIP process addresses both without managerial
authority: coordination through the immutable EIP stage framework,
cooperation through pseudonymous reputation. This is a new form in
Puranam et al.'s sense---the solution is not to reduce the problem to
an existing form but to deploy different mechanisms.

\noindent\textit{For technology standardization:} The comparison with
A2A reveals that high independent-contributor breadth ($\approx$44\% in both cases)
is compatible with informal institutional hierarchy: a single founding institution
(Google, 25.0\% of records) can dominate deliberation without formal authority~\cite{simcoe2012}. Openness is
a necessary but not sufficient condition for permissionless governance.
The distinguishing feature is the absence of a binding authority gate---in
A2A, the TSC gate; in SSOs, the membership gate~\cite{krechmer2006}.

\subsection{Limitations}

\textbf{Survivorship bias:} Both cases are successful protocols. Failed
governance attempts are not represented in the sample; the mechanisms
identified may be necessary but not sufficient for successful DAO governance.

\textbf{Comparison asymmetry:} ERC-8004 has 149 records; A2A has 5,272.
The structural difference in data volume reflects transparency norms but
makes direct statistical comparison difficult without normalization.
Proportional comparisons (argument type \%, stance \%) partially address this.

\textbf{Annotation reliability:} LLM annotation with MiniMax-M2.5 has not
been validated against a human-coded subsample. Institution labels were manually
verified for 109 of 626 authors (R07 investigation), correcting 40 labels;
517 remain LLM-inferred. [AUTHOR:
compute Cohen's $\kappa \geq 0.7$ for a random $N=50$ subsample before submission.]

\textbf{Off-platform governance:} A2A's critical decisions partially migrate
to Discord and TSC calls. PR \#1206 is the clearest case. The forum/GitHub
record for A2A is incomplete in ways the ERC-8004 record is not, which
biases the apparent governance transparency upward for ERC-8004.

\textbf{Scale confound:} ERC-8004 is a niche protocol at the time of
data collection; A2A has broader ecosystem adoption. Scale differences
may explain some behavioral differences independently of governance type.

\textbf{Single domain:} Both cases are AI agent protocols. Whether the
three mechanisms generalize to other DAO governance domains (DeFi protocol
governance, open-source infrastructure, creative commons) requires additional
research.


% =============================================================================
\section{Conclusion}
% =============================================================================

This paper examines how permissionless DAO governance achieves binding
consensus under conditions that transaction cost economics predicts require
hierarchy. Comparing ERC-8004 and Google A2A, we find that the DAO process
produces ratified consensus in 169 days with more genuine deliberative
engagement---higher Technical argument proportion, higher Oppose rate, more
equitable contribution distribution---while the corporate process exhibits
informal hierarchy, governance opacity at critical decision points, and
persistent institutional dominance.

Three mechanisms account for the DAO's coordination success. Smart contracts
provide the credible commitment that TCE assigns to authority. Pseudonymous
reputation provides the behavioral bonding that TCE assigns to employment.
Permissionless entry provides the competitive discipline that TCE assigns
to market exit. Together, they constitute what we term \textit{trustless
coordination}: governance without hierarchy, enforced by code, reputation,
and contestability.

These findings have practical implications for protocol design. Organizations
seeking to govern technical standards in high-specificity, high-uncertainty
domains face a choice: import hierarchy into an open process (the A2A path,
with attendant capture risk and opacity) or design for trustlessness from the
outset (the EIP path, with attendant coordination overhead and niche focus).
The EIP framework's success suggests that the perceived trade-off between
openness and decisiveness is at least partly false when the three mechanisms
are present.

[AUTHOR: add brief forward-looking sentence on implications for AI governance
infrastructure, connecting to the domain of ERC-8004's purpose---trustless
agent accountability---and how the governance process reflects the values it
encodes.]


% =============================================================================
\section*{Acknowledgment}
% =============================================================================

[AUTHOR: add funding acknowledgment if applicable. If none, delete this section.]


% =============================================================================
% REFERENCES
% All entries verified against data in paper/references.md.
% =============================================================================
\begin{thebibliography}{00}

% --- Transaction Cost Economics ---
\bibitem{williamson1975}
O.~E. Williamson, \textit{Markets and Hierarchies: Analysis and Antitrust
Implications}. New York: Free Press, 1975.

\bibitem{williamson1983}
O.~E. Williamson, ``Credible commitments: Using hostages to support
exchange,'' \textit{American Economic Review}, vol.~73, no.~4,
pp.~519--540, 1983.

\bibitem{williamson1985}
O.~E. Williamson, \textit{The Economic Institutions of Capitalism: Firms,
Markets, Relational Contracting}. New York: Free Press, 1985.

\bibitem{williamson1991}
O.~E. Williamson, ``Comparative economic organization: The analysis of
discrete structural alternatives,'' \textit{Administrative Science
Quarterly}, vol.~36, no.~2, pp.~269--296, 1991.

\bibitem{williamson1996}
O.~E. Williamson, \textit{The Mechanisms of Governance}. Oxford: Oxford
University Press, 1996.

\bibitem{klein1978}
B.~Klein, R.~G. Crawford, and A.~A. Alchian, ``Vertical integration,
appropriable rents, and the competitive contracting process,''
\textit{Journal of Law and Economics}, vol.~21, no.~2, pp.~297--326, 1978.

% --- Open-Source and Commons ---
\bibitem{omahony2007}
S.~O'Mahony and F.~Ferraro, ``The emergence of governance in an open source
community,'' \textit{Academy of Management Journal}, vol.~50, no.~5,
pp.~1079--1106, 2007.

\bibitem{shah2006}
S.~K. Shah, ``Motivation, governance, and the viability of hybrid forms in
open source software development,'' \textit{Management Science}, vol.~52,
no.~7, pp.~1000--1014, 2006.

\bibitem{west2003}
J.~West, ``How open is open enough? Melding proprietary and open source
platform strategies,'' \textit{Research Policy}, vol.~32, no.~7,
pp.~1259--1285, 2003.

\bibitem{simcoe2012}
T.~Simcoe, ``Standard setting committees: Consensus governance for shared
technology platforms,'' \textit{Management Science}, vol.~58, no.~8,
pp.~1509--1527, 2012.

\bibitem{krechmer2006}
K.~Krechmer, ``Open standards requirements,'' \textit{International Journal
of IT Standards and Standardization Research}, vol.~4, no.~1, pp.~43--61,
2006.

% --- DAO and Blockchain Governance ---
\bibitem{lumineau2021}
F.~Lumineau, W.~Wang, and O.~Schilke, ``Blockchain governance --- A new way
of organizing collaborations?'' \textit{Organization Science}, vol.~32, no.~2,
pp.~500--521, 2021.

\bibitem{defilippi2018}
P.~De~Filippi and A.~Wright, \textit{Blockchain and the Law: The Rule of
Code}. Cambridge, MA: Harvard University Press, 2018.

\bibitem{buterin2014}
V.~Buterin, ``A next-generation smart contract and decentralized application
platform,'' Ethereum White Paper, 2014. [Online]. Available:
\url{https://ethereum.org/en/whitepaper/}

\bibitem{nabben2023}
K.~Nabben, ``Is a `decentralized autonomous organization' a panopticon?
Algorithmic governance as creating and mitigating vulnerabilities in DAOs,''
\textit{Ledger}, vol.~8, pp.~1--23, 2023.

\bibitem{tan2022}
L.~Tan, S.~Teng, and Y.~Cao, ``Understanding governance and participation in
decentralized autonomous organizations,'' \textit{IEEE Transactions on
Engineering Management}, vol.~70, no.~8, pp.~2900--2913, 2022.

% --- New Forms of Organizing ---
\bibitem{puranam2014}
P.~Puranam, O.~Alexy, and M.~Reitzig, ``What's `new' about new forms of
organizing?'' \textit{Academy of Management Review}, vol.~39, no.~2,
pp.~162--180, 2014.

\bibitem{gulati1998}
R.~Gulati and H.~Singh, ``The architecture of cooperation: Managing
coordination costs and appropriation concerns in strategic alliances,''
\textit{Administrative Science Quarterly}, vol.~43, no.~4, pp.~781--814, 1998.

% --- Methods ---
\bibitem{eisenhardt1989}
K.~M. Eisenhardt, ``Building theories from case study research,''
\textit{Academy of Management Review}, vol.~14, no.~4, pp.~532--550, 1989.

\bibitem{yin2018}
R.~K. Yin, \textit{Case Study Research and Applications: Design and Methods},
6th ed. Thousand Oaks, CA: Sage, 2018.

% --- Platform and Standards ---
\bibitem{gawer2014}
A.~Gawer and M.~A. Cusumano, ``Industry platforms and ecosystem innovation,''
\textit{Journal of Product Innovation Management}, vol.~31, no.~3,
pp.~417--433, 2014.

\bibitem{north1990}
D.~C. North, \textit{Institutions, Institutional Change and Economic
Performance}. Cambridge: Cambridge University Press, 1990.

\end{thebibliography}

\vspace{12pt}

\end{document}
